\section{Related Work}
\label{sect:related_work}

\textbf{AI agent workflows.} Recent advances in LLM-based AI agents have introduced diverse workflows that combine language-based reasoning with external tool use. Single-agent frameworks (e.g., ReAct~\cite{react}, Reflexion~\cite{reflexion}, LATS~\cite{lats}, LLMCompiler~\cite{llmcompiler}) enhance decision-making through iterative reasoning, tool execution, and reflection. Multi-agent systems, such as CAMEL~\cite{camel} and AutoGen~\cite{autogen}, further extend these capabilities by structuring task execution, communication, and coordinated behaviors among multiple agents. 
While these workflows substantially improve their capabilities and behavioral flexibility, their system-level implications remain underexplored. This work provides the first comprehensive analysis of representative AI agents, offering insights into the efficiency and scalability of agentic systems.


\textbf{AI agent interfaces for tool-augmented reasoning.} In parallel with behavioral advancements in AI agents, recent efforts have focused on standardizing AI agent APIs and protocols to facilitate broader integration and deployment. OpenAI’s function-calling interface~\cite{openai_function_calling} defines a structured mechanism for API invocation, enabling agents to interact with tools in a verifiable and consistent manner. Anthropic’s Model-Context-Protocol (MCP)~\cite{mcp} further formalizes how agents manage context and interact with tools. Google’s Agent-to-Agent (A2A) protocol~\cite{agent2agent} complements these efforts by specifying a standard for multi-agent communication.
Although these contributions primarily standardize the interfaces and protocols for agent interaction, our work takes an orthogonal system-level perspective, uncovering the AI infrastructural challenges posed by agentic workloads under test-time scaling.

\textbf{System-level optimization of AI agents.}  
LLMCompiler~\cite{llmcompiler}, Alto~\cite{alto}, and Ayo~\cite{ayo} reduce inference latency by enabling pipelined and parallel execution across reasoning steps. Autellix~\cite{autellix} optimizes latency through queue-aware scheduling, while AI Metropolis~\cite{ai_metropolis} and Murakkab~\cite{murakkab} improve multi-agent coordination and resource isolation. 
While these works focus on optimizing specific components such as scheduling or execution flow, our study provides a broader characterization of infrastructural behaviors and efficiency trade-offs across diverse AI agents at scale.


\textbf{LLM inference optimization techniques.}
The AI community has only recently begun exploring agents, making it both unclear and highly challenging to determine the most effective methodology for applying various LLM-focused optimization techniques to agentic systems. To maintain clarity and generality in our analysis, this paper focuses on fundamental and widely adopted LLM inference optimizations that are readily available in existing AI serving frameworks~\cite{vllm,zheng2024sglang}. A comprehensive exploration of all the latest LLM optimizations in the literature for agentic systems is beyond the scope of this paper, so we provide a summary of recent LLM inference optimizations and discuss their applicability to agents below. 

In terms of KV cache management, \textit{hierarchical caching}~\cite{cachedattention,infinigen,multiturn} and \textit{non-prefix KV cache reuse}~\cite{cacheblend} approaches extend the naive prefix caching, enabling more efficient KV cache reuse. \textit{Token pruning}~\cite{dejavu,keyformer,attentionsink} or \textit{KV cache compression}~\cite{minicache,kvcachegen}, or model architectural improvement like grouped-query attention~\cite{ainslie2023gqa} and multi-head attention~\cite{liu2024deepseek} reduces the memory footprint of the KV cache, which will be especially helpful for agent workloads with long contexts.  Regarding decoding, \textit{Speculative decoding}~\cite{Leviathan2023SpecDecode} predicts multiple candidate tokens and validates in parallel to reduce decoding latency. In agents, speculative decoding can potentially become effective as their reasoning often generates predictable schema patterns (e.g., JSON structures or function arguments), which will increase the acceptance rate of speculative branches and improve overall decoding throughput. \textit{Prefill-decode disaggregation}~\cite{splitwise,distserve,dynamo} allows flexible and efficient resource allocation by decoupling the compute-intensive prefill phase from the memory-bound decode phase. For agents with long-context that incur substantial prefill computation load, disaggregation mitigates interference between prefill and decoding workloads, leading to more stable performance and improved overall efficiency.
